\documentclass[12pt,a4paper]{article}
\usepackage[UTF8]{ctex}
\usepackage{geometry}
\geometry{left=2.5cm,right=2.5cm,top=2.5cm,bottom=2.5cm}
\usepackage{graphicx}
\usepackage{booktabs}
\usepackage{multirow}
\usepackage{caption}
\usepackage{float}
\usepackage{amsmath}
\usepackage{hyperref}
\usepackage{xcolor}
\usepackage{enumitem}
\usepackage{fancyhdr}
\usepackage{titlesec}
\usepackage{array}

\pagestyle{fancy}
\fancyhf{}
\fancyhead[L]{\small 睡眠数据统计分析报告}
\fancyhead[R]{\small \thepage}
\renewcommand{\headrulewidth}{0.4pt}

\titleformat{\section}{\Large\bfseries}{\thesection}{1em}{}
\titleformat{\subsection}{\large\bfseries}{\thesubsection}{1em}{}

\hypersetup{colorlinks=true,linkcolor=blue!60!black}

\title{\textbf{两组睡眠障碍患者治疗效果\\对比分析报告}}
\author{数据分析项目}
\date{\today}

\begin{document}
\maketitle
\thispagestyle{empty}
\newpage
\tableofcontents
\newpage

\section{研究概况}

本研究纳入62例睡眠障碍患者，随机分为1组（对照组，n=31）和2组（观察组，n=31），比较两种治疗方案的疗效差异。主要观察指标：总睡眠时间（min）、睡眠效率（\%）、入睡后清醒时间（WASO，min）、睡眠潜伏期（min）。

疗效评定标准：根据四项指标综合改善率分为痊愈（$\geq 75\%$）、显效（$50\%\sim 75\%$）、有效（$25\%\sim 50\%$）、无效（$< 25\%$）。

\section{统计方法}

采用Shapiro-Wilk检验对各指标治疗前数据进行正态性检验，结果显示两组治疗前各指标均不满足正态分布假设（$P<0.05$），如表~\ref{tab:shapiro}所示。因此，描述性统计采用$M(P_{25}, P_{75})$表示，组内比较采用Wilcoxon符号秩检验，组间比较采用Mann-Whitney U检验，有效率比较采用Fisher精确检验。$\alpha=0.05$。

\begin{table}[H]
\centering
\caption{治疗前各指标Shapiro-Wilk正态性检验}
\label{tab:shapiro}
\begin{tabular}{llcc}
\toprule
\textbf{指标} & \textbf{组别} & \textbf{W值} & \textbf{P值} \\
\midrule
\multirow{2}{*}{睡眠时间(min)} & 1组 & 0.833 & $<0.001$ \\
 & 2组 & 0.917 & 0.003 \\
\midrule
\multirow{2}{*}{睡眠效率(\%)} & 1组 & 0.893 & $<0.001$ \\
 & 2组 & 0.853 & $<0.001$ \\
\midrule
\multirow{2}{*}{WASO(min)} & 1组 & 0.926 & 0.011 \\
 & 2组 & 0.886 & $<0.001$ \\
\midrule
\multirow{2}{*}{睡眠潜伏期(min)} & 1组 & 0.918 & 0.020 \\
 & 2组 & 0.909 & 0.012 \\
\bottomrule
\end{tabular}
\end{table}

\newpage
\section{基线资料比较}

治疗前两组各项睡眠指标比较差异均无统计学意义（$P>0.05$），说明两组基线具有可比性，见表~\ref{tab:baseline}。

\begin{table}[H]
\centering
\caption{两组患者治疗前睡眠指标基线比较 $[M(P_{25}, P_{75})]$}
\label{tab:baseline}
\begin{tabular}{lcccl}
\toprule
\textbf{指标} & \textbf{1组(n=31)} & \textbf{2组(n=31)} & \textbf{Z值} & \textbf{P值$^a$} \\
\midrule
睡眠时间(min) & $297.7(293.6, 344.6)$ & $348.1(294.1, 382.4)$ & 1.619 & $0.107$ \\
睡眠效率(\%) & $65.2(59.9, 72.2)$ & $73.0(65.8, 75.0)$ & 1.633 & $0.104$ \\
WASO(min) & $135.7(100.6, 159.9)$ & $98.7(72.7, 130.3)$ & 1.619 & $0.107$ \\
睡眠潜伏期(min) & $34.7(25.6, 40.1)$ & $28.9(24.9, 38.0)$ & 0.908 & $0.368$ \\
\bottomrule
\multicolumn{5}{l}{\small 注：$^a$Mann-Whitney U检验} \\
\end{tabular}
\end{table}

\section{治疗前后睡眠指标比较}

\subsection{组内比较}

两组治疗后各项睡眠指标均较治疗前显著改善（$P < 0.01$），见表~\ref{tab:pre_post}。

\begin{table}[H]
\centering
\caption{两组患者治疗前后睡眠指标比较 $[M(P_{25}, P_{75})]$}
\label{tab:pre_post}
\begin{tabular}{llccccc}
\toprule
\textbf{指标} & \textbf{组别} & \textbf{n} & \textbf{治疗前} & \textbf{治疗后} & \textbf{Z值} & \textbf{P值$^a$} \\
\midrule
\multirow{2}{*}{睡眠时间(min)} & 1组 & 31 & $297.7(293.6, 344.6)$ & $366.3(315.2, 401.8)$ & 2.548 & $0.010$ \\
 & 2组 & 31 & $348.1(294.1, 382.4)$ & $437.4(352.5, 476.2)$ & 2.724 & $0.005$ \\
\midrule
\multirow{2}{*}{睡眠效率(\%)} & 1组 & 31 & $65.2(59.9, 72.2)$ & $72.8(64.7, 84.5)$ & 2.587 & $0.009$ \\
 & 2组 & 31 & $73.0(65.8, 75.0)$ & $89.9(77.2, 96.8)$ & 3.351 & $0.001$ \\
\midrule
\multirow{2}{*}{WASO(min)} & 1组 & 31 & $135.7(100.6, 159.9)$ & $108.4(42.1, 142.9)$ & 2.587 & $0.009$ \\
 & 2组 & 31 & $98.7(72.7, 130.3)$ & $24.9(8.8, 53.4)$ & 3.410 & $0.001$ \\
\midrule
\multirow{2}{*}{睡眠潜伏期(min)} & 1组 & 31 & $34.7(25.6, 40.1)$ & $27.1(21.5, 35.2)$ & 2.567 & $0.009$ \\
 & 2组 & 31 & $28.9(24.9, 38.0)$ & $15.0(13.0, 24.6)$ & 3.096 & $0.001$ \\
\bottomrule
\multicolumn{7}{l}{\small 注：$^a$Wilcoxon符号秩检验} \\
\end{tabular}
\end{table}

\subsection{组间比较——治疗前后差值}

两组在四项指标的改善幅度上差异均有统计学意义（$P\leq 0.010$），见表~\ref{tab:diff}。

\begin{table}[H]
\centering
\caption{两组患者治疗前后差值组间比较 $[M(P_{25}, P_{75})]$}
\label{tab:diff}
\begin{tabular}{lcccl}
\toprule
\textbf{指标（差值）} & \textbf{1组(n=31)} & \textbf{2组(n=31)} & \textbf{Z值} & \textbf{P值$^a$} \\
\midrule
睡眠时间变化(min) & $38.1(0.6, 80.8)$ & $90.0(54.5, 130.5)$ & 2.626 & $0.009$ \\
睡眠效率变化(\%) & $5.6(-1.8, 19.3)$ & $22.5(5.2, 29.8)$ & 2.583 & $0.010$ \\
WASO变化(min) & $-36.4(-62.0, -1.1)$ & $-68.9(-98.4, -48.9)$ & 2.633 & $0.009$ \\
睡眠潜伏期变化(min) & $-7.4(-10.6, 1.5)$ & $-14.2(-19.0, -6.9)$ & 2.668 & $0.008$ \\
\bottomrule
\multicolumn{5}{l}{\small 注：$^a$Mann-Whitney U检验；差值=治疗后$-$治疗前} \\
\end{tabular}
\end{table}

\subsection{组间比较——治疗后}

治疗后2组在所有四项指标上均显著优于1组（$P<0.01$），见表~\ref{tab:post}。

\begin{table}[H]
\centering
\caption{两组患者治疗后睡眠指标组间比较 $[M(P_{25}, P_{75})]$}
\label{tab:post}
\begin{tabular}{lcccl}
\toprule
\textbf{指标} & \textbf{1组(n=31)} & \textbf{2组(n=31)} & \textbf{Z值} & \textbf{P值$^a$} \\
\midrule
睡眠时间(min) & $366.3(315.2, 401.8)$ & $437.4(352.5, 476.2)$ & 2.590 & $0.010$ \\
睡眠效率(\%) & $72.8(64.7, 84.5)$ & $89.9(77.2, 96.8)$ & 3.006 & $0.003$ \\
WASO(min) & $108.4(42.1, 142.9)$ & $24.9(8.8, 53.4)$ & 3.168 & $0.002$ \\
睡眠潜伏期(min) & $27.1(21.5, 35.2)$ & $15.0(13.0, 24.6)$ & 2.597 & $0.010$ \\
\bottomrule
\multicolumn{5}{l}{\small 注：$^a$Mann-Whitney U检验} \\
\end{tabular}
\end{table}

\subsection{治疗前后各指标对比可视化}

\begin{figure}[H]
\centering
\includegraphics[width=0.95\textwidth]{figures/fig1_pre_post_comparison.pdf}
\caption{两组患者治疗前后睡眠指标对比。*$P<0.05$，**$P<0.01$。}
\label{fig:pre_post}
\end{figure}

\newpage
\section{临床疗效评价}

\begin{table}[H]
\centering
\caption{两组患者临床疗效比较 [n(\%)]}
\label{tab:efficacy}
\begin{tabular}{lcccccc}
\toprule
\textbf{组别} & \textbf{n} & \textbf{痊愈} & \textbf{显效} & \textbf{有效} & \textbf{无效} & \textbf{总有效率} \\
\midrule
1组 & 31 & 0(0.0\%) & 0(0.0\%) & 14(45.2\%) & 17(54.8\%) & 45.2\% \\
2组 & 31 & 0(0.0\%) & 5(16.1\%) & 20(64.5\%) & 6(19.4\%) & 80.6\% \\
\bottomrule
\multicolumn{7}{l}{\small 注：总有效率比较采用Fisher精确检验，$P=0.008$} \\
\end{tabular}
\end{table}

2组总有效率（80.6\%）显著高于1组（45.2\%），差异有统计学意义（$P=0.008$）。

\begin{figure}[H]
\centering
\includegraphics[width=0.7\textwidth]{figures/fig3_efficacy_distribution.pdf}
\caption{两组患者临床疗效分布。TER=总有效率。**$P<0.01$。}
\label{fig:efficacy}
\end{figure}

\newpage
\section{多维度疗效分析}

\begin{figure}[H]
\centering
\includegraphics[width=0.9\textwidth]{figures/fig2_improvement_rate.pdf}
\caption{两组各指标改善率对比（中位数及四分位数间距）。**$P<0.01$。}
\label{fig:improve}
\end{figure}

\begin{figure}[H]
\centering
\includegraphics[width=0.65\textwidth]{figures/fig4_boxplot_improvement.pdf}
\caption{两组综合改善率分布。每个点代表一例患者。**$P<0.01$。}
\label{fig:boxplot}
\end{figure}

\begin{figure}[H]
\centering
\includegraphics[width=0.7\textwidth]{figures/fig5_radar_chart.pdf}
\caption{两组多维度改善率雷达图。2组在各维度上均优于1组。}
\label{fig:radar}
\end{figure}

\begin{figure}[H]
\centering
\includegraphics[width=0.95\textwidth]{figures/fig7_violin_plots.pdf}
\caption{各指标改善率小提琴图。}
\label{fig:violin}
\end{figure}

\begin{figure}[H]
\centering
\includegraphics[width=0.95\textwidth]{figures/fig6_correlation_heatmap.pdf}
\caption{两组改善率指标Spearman秩相关性热力图。}
\label{fig:corr}
\end{figure}

\newpage
\section{结论}

\begin{enumerate}[leftmargin=2em]
    \item 两组治疗前基线水平差异无统计学意义（$P>0.05$），具有可比性
    \item 两组治疗后各项睡眠指标均较治疗前显著改善（$P<0.01$）
    \item 2组在四项指标的改善幅度上均显著优于1组（$P\leq 0.010$）
    \item 2组总有效率（80.6\%）显著高于1组（45.2\%），$P=0.008$
    \item 综合以上结果，2组治疗方案疗效显著优于1组
\end{enumerate}

\end{document}
